\documentclass[../ppgc-tese.tex]{subfiles}
%\graphicspath{{\subfix{../images/}}}
\begin{document}
Although currently there are over eight billion living humans, we all share similar hardware, that lead to similar behavior in certain contexts. Almost all of us have a visceral understanding of basic emotions such as fear and joy. Although artificial models might behave very similarly to humans, they have an alien inner working that doesn't necessarily share the same inner workings. This creates a bias in the sense that we human believe that these models function and process experiences in the same way we do.

Although alignment is not necessary for understanding, understanding is paramount for proper alignment.

- Mechanistic interpretability

\begin{figure}
    \centering
    \caption{Model as a learned circuit}
    \begin{subfigure}[b]{0.5\textwidth}
        \centering
        \includegraphics[height=2cm]{a}
        \caption{Lorem ipsum}
    \end{subfigure}%
    ~ 
    \begin{subfigure}[b]{0.5\textwidth}
        \centering
        \includegraphics[height=2cm]{b}
        \caption{Lorem ipsum, lorem ipsum,Lorem ipsum, lorem ipsum,Lorem ipsum}
    \end{subfigure}
    \legend{Source: The authors, adapted from XXXX}
    \label{fig:model-circuits}
\end{figure}
\end{document}